\begin{definition}
Let $H'$ and $H$ be multihypergraphs on the same vertex type, with edge-index
types $F$ and $E$, respectively. We call $H'$ a sub-hypergraph of $H$ if
there exists an injective map $i:F\to E$ such that
$H'.\mathrm{edge}(f)=H.\mathrm{edge}(i(f))$ for every $f\in F$.
Thus inclusion preserves edge multiplicities: distinct child edge indices
use distinct ambient edge indices. Parallel ambient edges are allowed;
the ambient edge-to-vertex-set map need not be injective.
\end{definition}
